%\documentclass{beamer}
%\usepacakge[UTF8]{ctex}
%----------------------------------------------------------------------------------------
%    PACKAGES AND THEMES
%----------------------------------------------------------------------------------------

\documentclass[aspectratio=169,xcolor=dvipsnames]{beamer}
\usepackage[UTF8]{ctex}
\usetheme{SimplePlus}
\setbeamertemplate{footline}{}

\usepackage{hyperref}
\usepackage{graphicx} % Allows including images
\usepackage{booktabs} % Allows the use of \toprule, \midrule and \bottomrule in tables

\usepackage{mathpartir}
\usepackage{amssymb}  % 提供 \rightsquigarrow
\usepackage{graphicx} % 提供 \reflectbox

% 定义一个新的 \leftsquigarrow 命令
\newcommand{\leftsquigarrow}{%
  \mathrel{\reflectbox{$\rightsquigarrow$}}%
}

\newcommand\JEq{\mathtt{JEq}}
\newcommand\Ty{\mathtt{Ty}}
\newcommand\Ter{\mathtt{Ter}}
\newcommand\ttp{\mathtt{p}}
\newcommand\ttq{\mathtt{q}}
\newcommand\bfone{\mathbf{1}}
\newcommand\Set{\mathbf{Set}}
\newcommand\Id{\mathtt{Id}}
\newcommand\refl{\mathtt{refl}}
\newcommand\C{\mathbf{C}}
\newcommand\psC{\widehat{\mathbf{C}}}
\newcommand\extop{{\scalebox{0.4}[1]{\_}}.{\scalebox{0.4}[1]{\_}}}

\DeclareMathOperator\obj{obj}

%----------------------------------------------------------------------------------------
%    TITLE PAGE
%----------------------------------------------------------------------------------------

\title{族范畴(CwF)中的相等类型}
%\subtitle{Subtitle}

%\author{Pin-Yen Huang}

%\institute
%{
%    Department of Computer Science and Information Engineering \\
%    National Taiwan University % Your institution for the title page
%}
%\date{\today} % Date, can be changed to a custom date
\date{}

%----------------------------------------------------------------------------------------
%    PRESENTATION SLIDES
%----------------------------------------------------------------------------------------

\begin{document}

\begin{frame}
    % Print the title page as the first slide
    \titlepage
\end{frame}

\begin{frame}{背景}
\begin{itemize}
  \item 族范畴提供了解释依赖类型论中类型和项的基本框架
  \item 但是解释依赖和, 依赖积, 内涵相等类型需要更多的结构
\end{itemize}
\end{frame}

\begin{frame}{Id}
  $\Gamma\vdash A$, $\Gamma\vdash a:A$. CwF 支持内涵相等类型, 如果它支持以下结构
  \begin{itemize}
    \item 相等类型 \inferrule{\Gamma\vdash b:A}{\Gamma\vdash\Id_{A}(a,b)}
    \item 反身项 \inferrule{ }{\Gamma\vdash\refl_{A}(a):\Id_{A}(a,a)}
    \item 归纳原理\[ \inferrule{\Gamma.A.\Id_{A[\ttp]}(a[\ttp],\ttq)\vdash C \\ \Gamma\vdash c_{0}:C[a,\refl_{A}(a)] \\ \Gamma\vdash b:A \\ \Gamma\vdash u:\Id_{A}(a,b)}{\Gamma\vdash J_{A}(a,C,c_0,b,u):C[b,u]} \]
  \end{itemize}
  并且有以下等式(说的是这些构造和替换交换)
  \begin{align*}
  \Id_{A}(a,b)[\gamma] &= \Id_{A[\gamma]}(a[\gamma],b[\gamma]) \\
  \refl_{A}(a)[\gamma] &= \refl_{A[\gamma]}(a[\gamma]) \\
  J_{A}(a,C,c_0,b,u)[\gamma] &= J_{A[\gamma]}(a[\gamma],C[\gamma'],c_0[\gamma],b[\gamma],u[\gamma])
  \end{align*}
\end{frame}

\begin{frame}{$[a,b]$ 符号的解释}
  CwF 中有规则 \inferrule{\gamma:\Delta\to\Gamma \\ \Gamma\vdash A \\ \Delta\vdash a:A[\gamma]}{\langle\gamma,a\rangle:\Delta\to\Gamma.A}

  取特例
  \begin{itemize}
    \item \phantom{ }\inferrule{id:\Gamma\to\Gamma \\ \Gamma\vdash A \\ \Gamma\vdash a:A[id]=A}{\langle id,a\rangle:\Gamma\to\Gamma.A}
    \item \phantom{ }\inferrule{\langle id,a\rangle:\Gamma\to\Gamma.A \\ \Gamma.A\vdash B \\ \Gamma\vdash b:B[\langle id,a\rangle]}{\langle\langle id,a\rangle,b\rangle:\Gamma\to\Gamma.A.B}
  \end{itemize}
  把 $\langle\langle id,a\rangle,b\rangle$ 记为 $[a,b]$
\end{frame}

\begin{frame}{$  J_{A}(a,C,c_0,b,u)[\gamma] = J_{A[\gamma]}(a[\gamma],C[\gamma'],c_0[\gamma],b[\gamma],u[\gamma]$}
  \begin{itemize}
  \item \phantom{ }$\gamma:\Delta\to\Gamma$
  \item \phantom{ }\inferrule{\Delta.A[\gamma].\Id_{A[\gamma][\ttp]}(a[\gamma][\ttp],\ttq)\vdash C[\gamma'] \\ \Delta\vdash c_0[\gamma]:C[\gamma'][a[\gamma,\refl_{A[\gamma]}(a[\gamma])]] \\ \Delta\vdash b[\gamma]:A[\gamma] \\ \Delta\vdash u[\gamma]:\Id_{A[\gamma]}(a[\gamma],b[\gamma])}{\Delta\vdash J_{A[\gamma]}(a[\gamma],C[\gamma'],c_0[\gamma],b[\gamma],u[\gamma]):C[\gamma'][b[\gamma],u[\gamma]]}
  \item \phantom{ }$\gamma' :\Delta.A[\gamma].\Id_{A[\gamma][\ttp]}(a[\gamma][\ttp],\ttq) \to\Gamma.A.\Id_{A[\ttp]}(a[\ttp],\ttq)$
%  \item\phantom{ }\inferrule{\langle\gamma\circ\ttp,\ttq\rangle\circ\ttp:\Delta.A[\gamma].\Id_{A[\gamma][\ttp]}(a[\gamma][\ttp],\ttq) \to\Gamma.A \\ \Gamma.A\vdash\Id_{A[\ttp]}(a[\ttp],\ttq) \\ \Delta.A[\gamma].\Id_{A[\gamma][\ttp]}(a[\gamma][\ttp],\ttq) \vdash \ttq:\Id_{A[\ttp]}(a[\ttp],\ttq)[\langle\gamma\circ\ttp,\ttq\rangle\circ\ttp]}{}
  \item \phantom{ }\inferrule{\gamma\circ\ttp : \Delta.A[\gamma]\to\Delta\to \Gamma \\ \Gamma\vdash A \\ \Delta.A[\gamma]\vdash \ttq:A[\gamma\circ\ttp]=A[\gamma][\ttp]}{\langle\gamma\circ\ttp,\ttq\rangle:\Delta.A[\gamma]\to\Gamma.A}
  \end{itemize}
\end{frame}
\begin{frame}{续}
  \begin{itemize}
    \item \phantom{ }\inferrule{\langle\gamma\circ\ttp,\ttq\rangle\circ\ttp:\Delta.A[\gamma].\Id_{A[\gamma][\ttp]}(a[\gamma][\ttp],\ttq) \to \Delta.A[\gamma] \to \Gamma.A \\\\ \Gamma.A\vdash\Id_{A[\ttp]}(a[\ttp],\ttq) \\\\ \Delta.A[\gamma].\Id_{A[\gamma][\ttp]}(a[\gamma][\ttp],\ttq) \vdash \ttq:\Id_{A[\ttp]}(a[\ttp],\ttq)[\langle\gamma\circ\ttp,\ttq\rangle\circ\ttp]\\\\
    =\Id_{A[\ttp][\langle\gamma\circ\ttp,\ttq\rangle\circ\ttp]}(a[\ttp][\langle\gamma\circ\ttp,\ttq\rangle\circ\ttp],\ttq[\langle\gamma\circ\ttp,\ttq\rangle\circ\ttp])\\\\
    =\Id_{A[\gamma\circ\ttp\circ\ttp]}(a[\gamma\circ\ttp\circ\ttp],\ttq[\langle\gamma\circ\ttp\circ\ttp,\ttq[\ttp]\rangle)\\\\
    =\Id_{A[\gamma\circ\ttp][\ttp]}(a[\gamma\circ\ttp][\ttp],\ttq[\ttp])\\\\
    =\Id_{A[\gamma\circ\ttp]}(a[\gamma\circ\ttp],\ttq)[\ttp]  \qquad \text{$\leftsquigarrow$此处等式推理针对$q$的类型,说明 $q$ 的合理性}
    }{\langle\langle\gamma\circ\ttp,\ttq\rangle\circ\ttp,\ttq\rangle :\Delta.A[\gamma].\Id_{A[\gamma][\ttp]}(a[\gamma][\ttp],\ttq) \to\Gamma.A.\Id_{A[\ttp]}(a[\ttp],\ttq)}
  \end{itemize}
  (注: 由$\langle\_,\_\rangle$的唯一性, $\langle\gamma\circ\ttp,\ttq\rangle\circ\ttp=\langle\gamma\circ\ttp\circ\ttp,\ttq[\ttp]\rangle$)
  
  我们有如下定义
\[
\gamma'\triangleq\langle\langle\gamma\circ\ttp,\ttq\rangle\circ\ttp,\ttq\rangle
\]
\end{frame}

\begin{frame}{严格相等}
  \[
  J_{A}(a,C,c_{0},a,\refl_{A}(a)) = c_{0}
  \]
\end{frame}

\begin{frame}{弱相等}\[
  \inferrule{\Gamma.A.\Id_{A[\ttp]}(a[\ttp],\ttq)\vdash C \\ \Gamma\vdash c_{0}:C[a,\refl_{A}(a)]}{\Gamma\vdash\JEq_{A}(a,C,c_{0}):\Id_{C[a,\refl_{A}(a)]}(J_{A}(a,C,c_{0},a,\refl_{A}(a)),c_{0})} \]
  \pause
  满足等式\[\JEq_{A}(a,C,c_{0})[\gamma] = \JEq_{A[\gamma]}(a[\gamma],C[\gamma'],c_{0}[\gamma])\]
  \pause
  ($\gamma$, $\gamma'$ 如前)
\end{frame}

\begin{frame}{严格相等蕴含弱相等}
  严格相等意味着\[
    J_{A}(a,C,c_{0},a,\refl_{A}(a)) = c_{0}
  \]
  \pause
  此时有\[
  \inferrule{\Gamma.A.\Id_{A[\ttp]}(a[\ttp],\ttq)\vdash C \\ \Gamma\vdash c_{0}:C[a,\refl_{A}(a)]}{\Gamma\vdash\JEq_{A}(a,C,c_{0}):\Id_{C[a,\refl_{A}(a)]}(J_{A}(a,C,c_{0},a,\refl_{A}(a)),c_{0})=\Id_{C[a,\refl_{A}(a)]}(c_0,c_{0})} \]
  \pause
  只要定义
  \[
  \JEq_{A}(a,C,c_{0})\triangleq\refl_{C[a,\refl_{A}(a)]}(c_{0})
  \]
\end{frame}

\end{document}